% MWE: чи винен arydshln у "Extra }, or forgotten $"?
% Порядок пакетів — точно як у config/zkdl-template.cls (рядки 60-65).
\documentclass{article}
\usepackage[T1,T2A]{fontenc}
\usepackage[utf8]{inputenc}
\usepackage{amsmath,amssymb}
\usepackage{float}
\usepackage[dvipsnames,table]{xcolor}
\usepackage{tabularx}
\usepackage{blkarray}
\usepackage{arydshln}     % <<< ТУМБЛЕР: закоментуйте цей рядок і запустіть ще раз
\usepackage{makecell}
\begin{document}
\begin{table}[H]
    \centering
    \begin{tabularx}{\textwidth}{ccX}
        \Xhline{3\arrayrulewidth}
        \rowcolor{blue!30}\textbf{Section} & \textbf{Topic} & \textbf{Key Concepts} \\
        \Xhline{3\arrayrulewidth}
        \rowcolor{blue!10}\ref{section:number-theory} & Number Theory & Modular
        Arithmetic, Ring $\mathbb{Z}_n$ and field $\mathbb{Z}_n^{\times}$,
        Fermat's Little Theorem \\
        \hline
        \rowcolor{blue!20}\ref{section:abstract-algebra} & Abstract Algebra &
        Groups, Subgroups, Fields, Prime field $\mathbb{F}_p$, Isomorphisms,
        Automorphisms \\
        \hline
        \rowcolor{blue!10}\ref{section:polynomial-rings} & Polynomials & Divisibility, Lagrange Interpolation, Schwartz-Zippel Lemma \\
        \hline
        \rowcolor{blue!20}\ref{section:linear-algebra} & Linear Algebra Basics & Basic Operations over Vectors and Matrices \\
        \hline
        \rowcolor{blue!10}\ref{section:finite-fields} & Fields Extensions & Definition of $\mathbb{F}_{p^m}$, Field Multiplicative Subgroup, Algebraic Closure \\
        \Xhline{3\arrayrulewidth}
    \end{tabularx}
    \caption{Topics covered in Part I}
    \label{tab:essentialls}
\end{table}
\end{document}
